\section{Proposed Method}\label{sec:method}
% TODO: Present the new neuronal-network concept, intuition, formal definition and key equations.


% Example visualization: three rows of 8-bit sequences (MSB/LSB labeled) drawn with TikZ to illustrate the proposed encoding/neuronal input layout.
\begin{center}
\begin{tikzpicture}[scale=0.45,transform shape, box/.style={draw, minimum size=8mm, inner sep=1pt}, node distance=0.9cm]
    % sequence can be changed; this is a small "random" example
    \foreach \i [count=\xi from 0] in {1,0,1,1,0,0,1,0} {
        \node[box] (n\xi) at (\xi*0.95,0) {\i};
    }

    % labels above leftmost and rightmost boxes
    \node[above=2pt] at (n0.north) {LSB};
    \node[above=2pt] at (n7.north) {MSB};

    % add an extra box after a horizontal gap equal to 3 box widths
    % compute a random 0/1 for the extra box on the first row
    \pgfmathtruncatemacro{\rnda}{random(0,1)}
    \node[box] (nextra) at (11*0.95,0) {\rnda};

    % second row (random 0/1 sequence)
    \foreach \i [count=\xi from 0] in {0,1,0,0,1,1,0,1} {
        \node[box] (m\xi) at (\xi*0.95,-1.2) {\i};
    }
    % extra box after a gap of 3 boxes on the second row
    \pgfmathtruncatemacro{\rndb}{random(0,1)}
    \node[box] (mextra) at (11*0.95,-1.2) {\rndb};

    % three small dots between the last box (m7) and the extra rightmost box on the second row
    \node at (8.5,-1.2) {$\ldots$};
    \node at (8.5,-2.1) {$\ldots$};

    % centered triple dot between rows
    \node at (3.325,-2.1) {$\ldots$};

    % third row (another random 0/1 sequence)
    \foreach \i [count=\xi from 0] in {1,1,0,0,1,0,1,0} {
        \node[box] (p\xi) at (\xi*0.95,-3.0) {\i};
    }
    % extra box after a gap of 3 boxes on the third row
    \pgfmathtruncatemacro{\rndc}{random(0,1)}
    \node[box] (pextra) at (11*0.95,-3.0) {\rndc};

    % right-side vertical column (single column), placed relative to the extra box so it's visible
    \coordinate (colx) at ($(nextra.east)+(1.5,0)$);
    \pgfmathtruncatemacro{\colA}{random(0,1)}
    \node[box] (col1) at ($(colx)+(0,0)$) {\colA};
    \pgfmathtruncatemacro{\colB}{random(0,1)}
    \node[box] (col2) at ($(colx)+(0,-1.2)$) {\colB};
    % vertical dots indicating omitted rows
    \node at ($(colx)+(0,-2.1)$) {$\vdots$};
    \pgfmathtruncatemacro{\colC}{random(0,1)}
    \node[box] (col3) at ($(colx)+(0,-3.0)$) {\colC};
    % small vertical gap then more omitted rows
    \node at ($(colx)+(0,-4.2)$) {$\vdots$};
    \pgfmathtruncatemacro{\colD}{random(0,1)}
    \node[box] (col4) at ($(colx)+(0,-5.4)$) {\colD};
    \pgfmathtruncatemacro{\colE}{random(0,1)}
    \node[box] (col5) at ($(colx)+(0,-6.6)$) {\colE};
    \node at ($(colx)+(0,-7.5)$) {$\vdots$};
    \pgfmathtruncatemacro{\colF}{random(0,1)}
    \node[box] (col6) at ($(colx)+(0,-8.7)$) {\colF};

    % arrows from the first-row rightmost box (n7) to each of the extra rightmost boxes
    \draw[->,thick,>=latex] (n7.east) -- (nextra.west);
    \draw[->,thick,>=latex] (n7.east) -- (mextra.west);
    \draw[->,thick,>=latex] (n7.east) -- (pextra.west);
    % arrows from the second-row rightmost box (m7) to each of the extra rightmost boxes
    \draw[->,thick,>=latex] (p7.east) -- (nextra.west);
    \draw[->,thick,>=latex] (p7.east) -- (mextra.west);
    \draw[->,thick,>=latex] (p7.east) -- (pextra.west);
\end{tikzpicture}
\end{center}
